% 第6章 讨论：误差来源、失败案例、适用边界与关键影响因素（占位）
% 素材：experiments/dram-see-v3/blocker.md、diagnostic-report.md（预充失败复盘）；
%       experiments/dram-see-v4/source-field-audit-r1（粗网格空间失败保留）；
%       shared/literature/mentor-added-20261116.md（四类机制区分：瞬态收集/SEFI/位移损伤保持退化/BTI）；
%       缺口：A04 截面/错误率与外部验证（需 G01/G04 分母与导师实验数据 M01/M02）
\section{讨论}

\textbf{待补章节。} 本章计划讨论：数值误差 vs 判据误差 vs 物理机制的分层解释；失败案例（预充门槛、粗网格空间失败）的原因与边界；模型泛化到未见工况的能力；关键影响因素（尺寸、粒子、角度、温度等）的物理解释；与导师实验数据的对比（待交接）。
